\documentclass[12pt]{article}
\usepackage[utf8]{inputenc}
\usepackage[T1]{fontenc}
\usepackage{hyperref}
\title{Monpa, Tawang Language Resources for K-12 Curriculum Development}
\author{Compiled by Claude Sonnet 5 (claude-sonnet-5) for Andr\'as Kornai}
\date{}
\begin{document}
\maketitle

\section{Introduction}
Tawang Monpa (ISO 639-3 twm, Glottolog tawa1289) is also called Dakpa or Takpa. It is spoken in Tawang district, Arunachal Pradesh, India, and in northern Trashigang District, Bhutan. The identification notes list the family as ``Bookkeeping'', which is a Glottolog placeholder. Glottolog and Wikipedia place the language in East Bodish within Sino-Tibetan. Wikitongues rates it as vigorous on the EGIDS scale. Searches found scholarly descriptions but no school materials, primer, or dictionary. The list below is therefore short.

\section{Primary Educational Resources}
No primers, textbooks or school programmes for Tawang Monpa turned up in search results.

\section{Linguistic and Reference Materials}
\begin{itemize}
\item Tombleson, Anette Helgestad (SIL). Grammar sketch of Tawang Monpa, \emph{Himalayan Linguistics} 21(2), 2022. The search summary says it covers case and focus marking, tense, aspect and evidentiality, negation, pronouns and clause structure. \url{https://journalpub.escholarship.org/himalayanlinguistics/article/35127/galley/26158/download}
\item Related eScholarship item: \url{https://escholarship.org/uc/item/1cq303bg}
\item Lhama Wangchu (2002), \emph{An Introduction to the Tawang Monpa Language}, Government of Arunachal Pradesh, about 96 pages (Glottolog bibliographic record). \url{https://glottolog.org/resource/reference/id/40427}
\item Hyslop and Tshering (2008), preliminary notes on Dakpa (Tawang Monpa), Glottolog record. \url{https://glottolog.org/resource/reference/id/58824}
\item Hyslop and Tshering (2010), \emph{North East Indian Linguistics} 2, Glottolog record. \url{https://glottolog.org/resource/reference/id/679364}
\item Journal article on the copula system of Tawang Monpa (paywalled). \url{https://www.jbe-platform.com/content/journals/10.1075/ltba.17013.tom}
\item Another Benjamins catalogue entry returned for this language; I did not verify what it covers. \url{https://lwc1.benjamins.com/catalog/ltba.37.2.03bod}
\end{itemize}

\section{Religious and Community Resources}
None found in search results.

\section{Digital and Audio Resources}
Wikitongues language page (profile only; I did not verify whether it holds any recordings): \url{https://wikitongues.org/languages/twm/}

\section{Curriculum Development Context}
The reference description says the language is not known to be taught in schools. No dictionary or wordlist was found, and the writing system is not established in these results. A curriculum would have to be built from the Tombleson grammar and from new documentation work with the community.
\end{document}